\section{Historical unregularized encoders and their separate results}\label{sec:historical}

This part preserves the separate analyses of the earlier unregularized
algorithms. Its degree-range assumptions describe those historical
encoders only; the current regularized theorem has no logarithmic
lower bound. Their logarithmic-range hypotheses
are not hypotheses of \cref{thm:regularized}. Their positive-degree
fallback convention is replaced by the explicit floor in the primary
algorithm.

When all degrees are positive, form
\[
 S=D^{-1/2}AD^{-1/2}=I-L.
\]
Retain the degree eigenvector $D^{1/2}\one/\norm{D^{1/2}\one}_2$ and the
$K-1$ other eigenvectors with largest absolute eigenvalues in its orthogonal
complement. Let $U^\top U=I_K$ and $SU=U\Lambda$, where
$\Lambda=\diag(\mu_1,\ldots,\mu_K)$. On the proved event this is exactly
the informative $K$-dimensional space. The encoder supplies only
$(D,U,\Lambda)$ to the decoder. Put
\begin{equation}\label{eq:encoder-action}
 X=D^{-1/2}U,\qquad AX=D^{1/2}U\Lambda.
\end{equation}
Thus row actions on the retained space require no later edge lookup.
The normalized Laplacian eigenvalues contain the same information via
$\mu=1-\lambda_L$. Exact arithmetic and algebraic retained spaces define
the theorem; a fixed numerical rank threshold is not its definition.

The ordinary-SBM decoder uses the \emph{signed} reconstruction
\begin{equation}\label{eq:euclidean-decoder}
 Q=\operatorname{orth}(X),\quad \Pi_X=QQ^\top,\quad R_E=A\Pi_X.
\end{equation}
It fits free Bernoulli row profiles with the finite growing/replacement
schedule and mandatory final E/M step specified in
\cref{sec:sbm-proof}. Only probabilities are clipped; $R_E$ is not clipped.
Neither true labels nor a correct initial partition are inputs.

For DC-SBM we make the explicit degree-metric correction
\begin{equation}\label{eq:metric-decoder}
 H_D=D^{-1/2}UU^\top D^{1/2},\qquad
 R_D=AH_D=D^{1/2}U\Lambda U^\top D^{1/2}.
\end{equation}
With $h_i=\norm{U_{i\cdot}}_2^2$ and
$s_i=(R_D)_{ii}/(1-h_i)$, its hollow version is
\begin{equation}\label{eq:hollow-decoder}
 R^\circ=R_D+\diag(s)(H_D-I).
\end{equation}
It satisfies
$\diag(R^\circ)=0$, $R^\circ\one=D\one$ and $R^\circ X=AX$.
It need not be symmetric after hollowing; a symmetric block parameter is
fitted with the ordered-row working criterion divided by two.

For the initial partition, normalize every row of $U$ to unit length.
Repeatedly select a remaining row having the largest number of remaining
rows within distance $1/8$, and remove rows within distance $1/4$ of the
selected row. After $K$ selections assign every row to the closest selected
row. The separation argument in \cref{sec:dcsbm-proof} proves its
$O(1/d)$ error; this is a deterministic procedure with $O(Kn^2)$ distance
comparisons, not an assumed global k-means solution.
Normalize degree estimates to mean one within each fitted group, fit
the symmetric block matrix by the exact Bernoulli working criterion, and
perform at least one exact Bernoulli profile-LRT assignment using
$R^\circ$. One pass suffices. The same theorem permits a finite number of
the stated hard refits, including adaptive stopping, by the invariant-basin
argument proved below.

Both algorithms return a fixed partition, or their initial partition,
if a prescribed operation is undefined. This includes zero degrees,
insufficient retained rank, a hollow denominator zero, an empty initial
group, or an inadmissible parameter fit. The required operations are all
well-defined on the good event, so this convention costs only its actual
failure probability.

\begin{theorem}[Earlier ordinary free-profile decoder]\label{thm:sbm}
Fix $K\ge2$, $\pi_->0$ and $\kappa>0$. Suppose
\[
 n_a\ge\pi_-n,\qquad
 p/2\le P_{ab}\le p\le1/4,\qquad
 \sigma_{\min}(P)\ge\kappa p,\qquad d=np\ge5\log n.
\]
Use \eqref{eq:euclidean-decoder} and the finite growing/replacement
likelihood decoder, with probability floor
$\varepsilon=c_0\max(1,n^{-1}\sum_iD_i)/n$ for a fixed sufficiently small
$c_0>0$ (for example any $c_0<1/16$ on this class).
There are constants $C,n_0$, depending only on the displayed fixed
structural constants and $c_0$, such that for $n\ge n_0$,
\begin{equation}\label{eq:sbm-risk}
 \E\ell(\widehat z_E,z)\le C R_{\mathrm{SBM},*}(n).
\end{equation}
The algorithm uses neither $P$, $p$, nor the true partition.
This includes fixed arbitrary $K$, unbalanced group sizes of comparable
order, and general symmetric full-rank block probability matrices in the
displayed range. It does not require all within-group probabilities, or all
between-group probabilities, to be identical.
\end{theorem}

\begin{theorem}[Earlier unregularized DC profile decoder]
\label{thm:dc}
Assume \eqref{eq:mean-model}--\eqref{eq:structural}. Fix $C_1<\infty$ and
suppose
\begin{equation}\label{eq:dc-range}
 C_0\log n\le d=np,\qquad np^2\le C_1,
\end{equation}
where $C_0$ is sufficiently large depending only on the fixed structural
bounds and $C_1$. The degree-metric decoder
\eqref{eq:metric-decoder}--\eqref{eq:hollow-decoder}, with the deterministic
initialization and estimated-parameter profile procedure just specified,
satisfies
\begin{equation}\label{eq:dc-risk}
 \E\ell(\widehat z_D,z)\le C R_{\mathrm{prof},*}(n)
\end{equation}
for all sufficiently large $n$, with $C$ independent of $n$.
The statement is uniform over the displayed compact parameter class.
No correct grouping, true degree parameters, true block matrix, sample
splitting, or leave-one-out eigensystem is an input to the algorithm.
\end{theorem}

\begin{corollary}[The known-target-degree oracle on equal exposures]
\label{cor:known}
Under \cref{thm:dc}, suppose additionally that the actual finite-$n$ block
exposures satisfy
\begin{equation}\label{eq:equal-exposure}
 \sum_b(n_b/n)C_{ab}=S\quad\text{for all }a.
\end{equation}
Then
\[
 R_{\mathrm{prof},*}\asymp R_{\theta,*},\qquad
 \E\ell(\widehat z_D,z)\le C R_{\theta,*}.
\]
Thus there are nontrivial fixed-$K$ SBM and DC-SBM classes on which the
spectral-only fitted likelihood procedure attains a constant multiple of
the original known-parameter local LRT risk. The ordinary free-profile
decoder, with a sufficiently small floor coefficient for its structural
parameter class, also satisfies \eqref{eq:sbm-risk} for general positive
full-rank $C$ under \eqref{eq:dc-range} and \eqref{eq:equal-exposure},
without the entry-ratio-two restriction of \cref{thm:sbm}.
\end{corollary}

Exact equality in \eqref{eq:equal-exposure} can be relaxed to exposure
differences whose contribution to all likelihood contrasts is uniformly
$O(1)$. For example fixed $C\pi=S\one$ with
$n_a=n\pi_a+O(1)$ suffices. Equality of limiting exposures with no rate
is not by itself sufficient. The general known-$\theta_i$ claim without
an additional restriction is false: \cref{thm:imposs} gives an explicit
diverging risk ratio for every full-graph algorithm.

\begin{remark}[What is and is not constant]
The factor may depend on signal conditioning, group proportions, the
allowed degree-parameter range and $C_1$. It is independent of $n$ on
the stated class. We do not assert a small numerical constant, a limit of
the risk ratio, or uniformity when $\sigma_{\min}(C)\to0$. The latter
weak-signal extension needs separate quantitative estimates.
\end{remark}

